\section{Introduction}\label{sec:introduction}

An integral polymatroid describes nonnegative integer allocations subject
to submodular capacity constraints. Intersecting two such systems is a
natural common generalization of matroid intersection and capacitated
transportation problems. Exact feasibility and optimization are accessible
through rank-value oracles. Counting the feasible integer allocations
poses additional difficulties: a large capacity permits exponentially
many coordinate values, and replacing it by labelled copies changes both
the size of the instance and the weight assigned to each allocation.

This paper gives an approximation algorithm whose cost depends
polynomially on the binary encoding of the capacities. The proof builds
on the common-matroid-base counting framework and its two-lift extension
for cell-bounded contingency tables~\cite{Parent,CB}. The new reduction
keeps only part totals with polynomially many possible values as discrete variables. It treats
the remaining degrees of freedom by weighted volume approximation and
proves the discrete quadratic inequalities needed by the counting
framework after this integration.

\subsection{The problem and the result}

Let $E=\{1,\ldots,n\}$. For $i\in\{1,2\}$, let
$r_i:2^E\to\Z_{\ge0}$ be normalized, nondecreasing, and submodular:
\[
 r_i(\varnothing)=0,\qquad
 r_i(A)\le r_i(B)\quad(A\subseteq B),\qquad
 r_i(A)+r_i(B)\ge r_i(A\cup B)+r_i(A\cap B).
\]
Assume their total ranks agree, and write $R=r_1(E)=r_2(E)$.
For $x\in\R^E$ and $A\subseteq E$, put $x(A)=\sum_{j\in A}x_j$.
The set of common integer bases is
\begin{equation}\label{eq:counting-problem}
 \Omega(r_1,r_2)=
 \left\{x\in\Z_{\ge0}^E:
 x(E)=R,\quad x(A)\le r_i(A)
 \text{ for every }A\subseteq E,\ i\in\{1,2\}\right\}.
\end{equation}
Write $Z=\abs{\Omega(r_1,r_2)}$. Every vector has weight one.
A representation of a vector by labelled
copies is not another object to be counted.

The input supplies $R$ in binary and exact value oracles for $r_1,r_2$.
A query is an $n$-bit subset indicator; its answer is the binary integer
$r_i(A)$, between zero and $R$. No decomposition, explicit inequality
list, sampler, or counting oracle is part of the input. Let $\ell$ be
the binary encoding length of $n,R,\eps,\delta$, where the rational
numbers $0<\eps,\delta<1$ are given by their numerators and denominators.
Oracle computation itself is excluded from the bit-operation count, but
writing queries and reading answers are included.

\begin{theorem}\label{thm:main}
There is a single uniform randomized algorithm which, for every input
described above, outputs a nonnegative rational $\widehat Z$ satisfying
\[
 \PP\bigl((1-\eps)Z\le\widehat Z\le(1+\eps)Z\bigr)
 \ge1-\delta.
\]
If $Z=0$, the algorithm outputs zero on every execution. On every
execution, both its number of oracle calls and its number of other bit
operations are bounded by one fixed polynomial in
\[
 n,\qquad \ell,\qquad \eps^{-1},\qquad \log(\delta^{-1}).
\]
The bit-operation bound includes rational arithmetic, generation of
independent unbiased random bits, and the encoding of the output.
\end{theorem}

The theorem includes lower-dimensional feasible sets, arbitrary pairs
of rank functions, and binary ranks of unrestricted magnitude. When
$R=0$, there is one common base, the zero vector; this includes the empty
vector when $n=0$. The algorithm's polynomial has large fixed exponents.
The issue here is dependence on the input encoding, not an optimized
running time.

\subsection{Context and prior work}

The polyhedral theory begins with Edmonds' work on submodular functions,
matroids, and polymatroid intersection~\cite{Edmonds1970}. It explains
why the intersection considered here is integral and why optimization
admits an oracle formulation. The exact algorithms used below come from
submodular minimization~\cite{IFF2001} and the rational
separation--optimization theory of Gr{\"o}tschel, Lov{\'a}sz, and
Schrijver~\cite{GLS1988}. The representation of an integer polymatroid by
a matroid on groups of interchangeable elements also has a classical
history, in Helgason's theory of hypermatroids~\cite{Helgason1974}.
This representation preserves the feasible occupancy vectors, but its
explicit size and the number of labelled realizations of a vector matter
for counting. The distinction becomes essential when capacities are
encoded in binary.

Contingency tables provide a concrete instance of this distinction.
Such a table is a nonnegative integer matrix with prescribed row and
column sums; cell bounds impose additional upper bounds on its entries.
Dyer, Kannan, and Mount obtained approximate counting and sampling
algorithms under lower bounds on the margins, using the relation between
lattice counts and transportation-polytope volumes~\cite{DKM1997}.
Cryan and Dyer gave an FPRAS with arbitrary margins when the number of
rows is fixed~\cite{CD2003}, and Cryan, Dyer, and Randall extended that
fixed-row result to cell-bounded tables~\cite{CDR2010}. For fixed-row
tables without cell bounds, Gopalan et al.\ subsequently gave a
deterministic FPTAS~\cite[Theorem~1.3]{GKM2011}. These results
already address numerical data through their binary encodings, within
their respective restrictions. For tables with $0/1$ entries,
Bez{\'a}kov{\'a}, Bhatnagar, and Vigoda gave an annealing-based FPRAS
for arbitrary margins~\cite{BBV2007}. The companion manuscript~\cite{CB} treats
cell-bounded tables with both dimensions variable. Its paired encoding
and two factorial corrections are the immediate counting tools used
here; the present reduction must additionally handle arbitrary
submodular rank constraints available only through value queries.

Another line of work develops the concavity behind matroid counting.
Adiprasito, Huh, and Katz proved the Hodge--Riemann relations for matroid
Chow rings and deduced log concavity of characteristic-polynomial
coefficients~\cite{AHK2018}. Anari, Liu, Oveis Gharan, and Vinzant used
log-concave polynomials and high-dimensional walks to obtain an
independence-oracle FPRAS for the bases of a single arbitrary
matroid~\cite{ALOV2024}. Br{\"a}nd{\'e}n and Huh's theory of Lorentzian
polynomials includes matroid basis polynomials and factorial-normalized
generating polynomials of integer polymatroid bases~\cite{BH2020}.
These statements concern an individual matroid or polymatroid; taking
common bases need not preserve the exchange property of its support.
Moreover, the factorial normalization weights integer bases nonuniformly.
They therefore do not directly yield the uniform common-base count
considered here. Our proof verifies the needed quadratic inequalities
for its particular auxiliary weights.

For common bases of two rank-$r$ matroids, Anari, Oveis Gharan, and
Vinzant obtained a deterministic polynomial-time $2^{O(r)}$-approximation
and a randomized relative approximation with running time
$2^{O(r)}\operatorname{poly}(n,\varepsilon^{-1},\log\delta^{-1})$
\cite[Theorem~1.3 and Corollary~1.4]{AOV2021}.
The companion manuscript~\cite{Parent} gives an FPRAS for common bases
of two arbitrary matroids on an explicitly listed ground set.
It uses positive rank-deficiency weights, transport of a single defect,
variance bounds for specified observables, and mixing of a trace chain.
The cell-bounded companion~\cite{CB} develops the two complementary
factorial lifts that cancel labelled multiplicities on the target
states. These constructions are related to the perfect- and near-perfect
matching ensemble and learned defect weights in the permanent FPRAS of
Jerrum, Sinclair, and Vigoda~\cite{JSV2004}. Restricted Poincar{\'e}
inequalities for the hole-pattern observables identifying unmatched
vertices also appear in the recent permanent algorithm of Chen,
Vigoda, and Yang~\cite{CVY2026}. The formal transport
and counting inputs to this paper are the specific statements from
the two companions; our contribution is the reduction to an instance
whose explicit size is polynomial in the original binary input.

The continuous part uses two classical results with different roles.
Pr{\'e}kopa's marginal theorem preserves log concavity when free
coordinates are integrated out~\cite[Theorem~6]{Prekopa1973}.
The randomized convex-volume algorithm of Dyer, Frieze, and
Kannan~\cite{DFK1991}, in the rational oracle formulation described by
Dyer and Frieze~\cite{DyerFrieze1991}, supplies approximate evaluations of
the resulting weights. Our rational-center construction uses determinant
replacement, closely related to the barycentric-spanner construction of
Awerbuch and Kleinberg~\cite[Proposition~2.4]{AK2004}; we give the affine
simplex argument needed here. These continuous tools supply the weights
and their evaluations, while the discrete signature argument establishes
the additional inequalities required by the counting chain.

\subsection{Proof strategy}

The first obstacle is the scale of the coordinate ranges. The expansion
matroid of an integral polymatroid can represent each unit of capacity by
a separate element, but a direct construction can have size proportional
to $nR$. In addition, an occupancy vector has a product of binomial
coefficients worth of labelled realizations. Applying a common-base
counter to that expansion would therefore address neither the required
runtime nor the required weights.

We begin with an exact feasibility test and a rational center of the
common base polytope. Near-tight rank inequalities at this center define
two partitions of $E$. Their part totals will be the discrete variables
retained by the algorithm. A coarser partition identifies directions in
which pairs of bases can be shifted without significantly changing their
count. Directional log concavity makes this comparison quantitative.
Within suitable fixed-total slices, a small homothety then shows that
enlarging the original base polytopes in the remaining directions changes
the count by only a small relative error.

For the retained part totals, an incidence graph joins a part of the
first partition to a part of the second whenever they share a coordinate.
Choosing spanning trees expresses all other coordinates through free
variables and the part totals. The retained integer ranges have polynomial
cardinality. Large free ranges are integrated out, with soft penalties
for rank violations. The preceding comparison proves that the resulting
weighted count approximates $Z$.

Integration preserves ordinary log concavity, but that property alone
does not give the discrete quadratic inequalities used by the counting
algorithm. We convolve the integrated weights with a smooth compactly
supported kernel and add a small strictly concave quadratic term.
Derivative estimates bound the errors between discrete second differences
and the continuous Hessian. The added curvature dominates these errors
at every annealing power, giving both required quadratic tests throughout
the counting algorithm.

Finally, each bounded part-total variable is represented by a polynomial
number of paired labels. Two factorially corrected lifts agree on
transversals, where their multiplicities cancel exactly, and satisfy the
complementary signature conditions of the paired counting framework.
We state the imported transport and trace estimates precisely, verify
their hypotheses, and give the capped annealing algorithm. A separate
convex-volume construction evaluates its weights with bounded rational
computation. This last distinction is essential: all large expansions and
slice subdivisions used for comparison are proof devices, whereas every
instance actually sampled has polynomial size.

Section~\ref{sec:polyhedra} supplies the oracle and polyhedral tools.
Section~\ref{sec:signatures} proves the integer signature and shift-count
statements. Section~\ref{sec:geometry} establishes the geometric
comparison, and Section~\ref{sec:ports} constructs the regularized
polynomial-range model. Section~\ref{sec:evaluation} evaluates its
weights. Section~\ref{sec:counting} performs the paired counting and
proves Theorem~\ref{thm:main}.
